\chapter{Implementation}
\label{cap:implementation}

This chapter describes how the design of Chapter~\ref{cap:design} was built.
It states the technology stack and its rationale, the layout of the
codebase, and then walks through the principal modules --- the talent pool,
the CV parsing pipeline, jobs and matching, the AI interviewer integration,
notifications, analytics, authentication and observability --- before
covering the continuous-integration pipeline and a small set of
representative code excerpts.

%----------------------------------------------------------------------------------------

\section{Technology Stack}
\label{sec:impl-stack}

Every technology was weighed against five criteria: fitness for purpose,
developer experience, type-safety integration with TypeScript, compatibility
with Vercel's serverless runtime, and the ability to run on free or student
tiers. The stack deliberately avoids heavyweight enterprise infrastructure
(Kubernetes, message queues, microservices) in favour of a right-sized
single application. Table~\ref{tab:tech-stack} summarises the choices.

\begin{table}[h]
\centering
\caption[Technology stack]{Principal technologies and their roles.}
\label{tab:tech-stack}
\small
\begin{tabular}{p{0.22\linewidth}p{0.34\linewidth}p{0.32\linewidth}}
\toprule
\textbf{Concern} & \textbf{Technology} & \textbf{Rationale} \\
\midrule
Framework & Next.js 16 (App Router), TypeScript & One codebase for UI and API; file-based routing; Vercel-native \\
UI & shadcn/ui + Tailwind CSS & Accessible primitives without a heavy framework \\
Authentication & Supabase Auth (Google OAuth) & Managed identity; role claim in the JWT \\
Database & Supabase PostgreSQL (no ORM) & Relational domain; direct typed client \\
Storage & Supabase Storage & Uploaded CV files \\
LLM & Groq Llama 3.3 70B (primary), OpenAI GPT-4o (fallback) & Free-tier primary with a reliability fallback \\
Validation & Zod & Runtime validation of all external/LLM input \\
Email & Resend & Assessment invitations and outreach \\
Charts & Recharts & Analytics dashboards \\
Testing & Vitest + V8 coverage & Fast, ESM-native unit testing \\
Interviewer & Python FastAPI (sidecar) & Separate process for conversational AI \\
CI & GitHub Actions & Type-check + test on every change \\
\bottomrule
\end{tabular}
\end{table}

The single most consequential choice was using \textbf{one framework for
both frontend and backend}, which removes a separate server, shared-type
generation and CORS handling. The alternatives --- a separate React/Express
backend, Angular/NestJS, Vue/Fastify or Remix --- were rejected as either
duplicative or a poor fit for the component ecosystem and the deployment
target.

%----------------------------------------------------------------------------------------

\section{Repository Layout}
\label{sec:impl-layout}

The repository is organised so that the onion layers are visible in the
folder structure. Under \texttt{src/}:

\begin{description}
	\item[\texttt{server/domain/}] value objects, pure domain services
	(scoring, job fit) and the port interfaces --- zero external
	dependencies.
	\item[\texttt{server/application/}] use cases (one file per module),
	DTOs and the canonical application-error types.
	\item[\texttt{server/infrastructure/}] Supabase repositories, the AI
	parsing services, storage, email and logging --- the only code that
	imports external SDKs.
	\item[\texttt{app/}] Next.js App Router: pages and the API route
	handlers under \texttt{app/api/}.
	\item[\texttt{client/}] React components, including the shadcn/ui
	primitives and the analytics widgets.
	\item[\texttt{lib/}] cross-cutting helpers, including the Supabase
	browser/server clients and the \texttt{require-hr} authorisation
	helper.
\end{description}

Outside \texttt{src/}, \texttt{tests/} holds the Vitest suite,
\texttt{supabase/migrations/} the canonical schema, and \texttt{scripts/}
the operational backfill utilities.

%----------------------------------------------------------------------------------------

\section{Talent Pool Module}
\label{sec:impl-pool}

The talent pool is the system of record for candidates. It supports creation
and editing of candidate records, full-text-style search and filtering (by
status, country, score range, language and source type), sorting, tagging and
free-text recruiter notes. A use case in the Application layer orchestrates
each operation and delegates persistence to a candidate repository through
its port, so the route handlers stay thin and the business rules stay
testable. Search terms supplied by the user are escaped before they reach the
PostgREST \texttt{.or()} filter (Section~\ref{sec:impl-code}), closing the
injection vector found during the audit.

%----------------------------------------------------------------------------------------

\section{CV Parsing Pipeline}
\label{sec:impl-cv}

CV ingestion is the most complex feature and is implemented as a nine-stage
pipeline with two entry points --- candidate self-upload (a single file) and
HR bulk upload (many files or a ZIP archive). The stages are: (1) validate the
file's MIME type and size; (2) store the original in Supabase Storage; (3)
extract text (\texttt{unpdf} for PDF, \texttt{mammoth} for DOCX, direct
decoding for plain text); (4) validate that enough text was recovered; (5)
send the text to the LLM with a structured-output prompt; (6) validate the
model's JSON against a strict Zod schema, retrying once on failure; (7) run
deduplication against existing candidates; (8) upsert the candidate and its
experiences, education, languages and skills; and (9) compute the CV score.

The prompt instructs the model to return a single JSON object with personal
information, experiences (each with a job title, company, dates, current-role
flag and a field-of-work tag), education, self-declared language levels and
categorised skills, using the provider's JSON mode. Two failure modes
dominate and are handled explicitly: \textbf{schema drift}, where the model
returns slightly malformed or invented fields --- absorbed by the tolerant
Zod schema and the single retry --- and \textbf{provider rate limits or
outages}, absorbed by retrying the same request against the OpenAI fallback.
A telemetry row is written per extraction for observability. For bulk
uploads, a \texttt{parsing\_jobs} row tracks parsed, failed and total counts
so that HR can monitor progress by polling.

%----------------------------------------------------------------------------------------

\section{Jobs and Matching}
\label{sec:impl-jobs}

Jobs are created and edited by HR or ingested by a scraper that collects
real listings from the adidas careers portal (used to populate a realistic
demonstration dataset). When an HR user first opens a job's ranking screen,
its description is parsed lazily into the versioned requirements schema and
cached. The orchestrator \texttt{matchCandidatesToJob} then loads the
candidate pool with their experiences, languages, education and skills,
builds the per-field experience vector each candidate needs, runs the pure
\texttt{computeJobFit} for every candidate, and persists the top results into
the \texttt{job\_matches} cache. The \texttt{match-candidates} screen renders
the ranking with the per-criterion breakdown, and a per-job
\texttt{job\_shortlists} table records the recruiter's working pick list,
snapshotting the fit score at the moment a candidate is added so that later
re-ranking does not rewrite history.

A subtle implementation point lives in the bridge that converts a persisted
database row into the matcher's input type. An early version read the wrong
property for the job title, which silently emptied the evidence text used for
title-based seniority matching; the defect was invisible because the pure
matcher itself was correct. It is now fixed and locked by the
\texttt{job-matching-bridge} tests of Chapter~\ref{cap:validation}, and it is
the concrete basis for the testing lesson in Chapter~\ref{cap:conclusions}.

%----------------------------------------------------------------------------------------

\section{AI Interviewer Integration}
\label{sec:impl-interviewer}

The conversational interviewer itself is the second team member's
contribution; this section documents only the integration surface owned by
the Next.js application. On the application side the flow is: an API route
creates an \texttt{interview\_sessions} row and issues a short-lived
HMAC-SHA256 token (storing only its hash); the browser opens an interview
popup that streams audio and chat; transcript turns are recorded through a
realtime turn endpoint; proctoring events are logged; and, on completion, the
structured evaluation returned by the sidecar is saved as the assessment
result. The selected skill (for the technical mode) is persisted on the
session and enforced server-side so that questions cannot drift off the
chosen topic. The token-signing utility resolves its secret at call time
rather than at module load, so that a local build can complete without the
deployment-only secret while still failing fast at request time in
production.

%----------------------------------------------------------------------------------------

\section{Notifications and Campaigns}
\label{sec:impl-notifications}

The notification system serves both roles. It records pipeline events
(status changes, contact emails, promotional outreach) as typed
notifications, each carrying a metadata payload and the HR sender where
applicable. A per-candidate interaction-history panel reconstructs the full
timeline --- status changes, emails (with expandable bodies) and campaign
membership --- by joining notifications to their originating campaign or job.
Promotional campaigns are authored in a rich-text editor and fan out to
preference-aware recipients, respecting each candidate's field and country
notification preferences so that irrelevant outreach is suppressed.

%----------------------------------------------------------------------------------------

\section{Analytics}
\label{sec:impl-analytics}

The analytics dashboard combines seven built-in charts (pipeline funnel,
status distribution, top skills and languages, trends and so on) with a
constrained custom-widget builder. Rather than exposing free-form SQL, the
builder is driven by a Zod-validated catalogue that whitelists the available
metrics, the dimensions allowed per metric, the chart types allowed per
dimension family, and the permitted filter keys; a strict widget-specification
schema rejects any unknown key. A hand-written query service runs each
validated specification, and a single universal renderer turns the resulting
\texttt{\{label, value\}} series into a bar, line, area, pie or stat chart.
Saved widgets are stored per user. This design gives HR analytical freedom
while making it impossible to express an injected or malformed query
(Chapter~\ref{cap:validation}).

%----------------------------------------------------------------------------------------

\section{Authentication, Roles and Middleware}
\label{sec:impl-auth}

Authentication uses Supabase Google OAuth. On first sign-in, an auth callback
route assigns the user's role into \texttt{app\_metadata.role} using the
service-role key, making it immutable from the client. The Next.js middleware
runs on every request: it refreshes the session, returns \texttt{401} for
unauthenticated API callers, returns \texttt{403} for non-HR callers on the
HR-only API prefixes, and protects the dashboard routes. Keeping this logic
in one middleware --- rather than repeating it in every route handler --- means
a new route is gated by default and the check cannot be forgotten. Inside
route handlers that need it, a small \texttt{require-hr} helper performs the
same role check with the revocation-aware \texttt{getUser()} call.

%----------------------------------------------------------------------------------------

\section{Logging and Observability}
\label{sec:impl-observability}

Structured logging is provided by an in-house \texttt{createLogger(scope)}
factory that replaced raw \texttt{console} calls across the API routes. It
emits scoped, structured records and redacts personally identifiable
information --- e-mail addresses and tokens are masked --- so that logs
shipped to the hosting platform do not leak candidate data, a property pinned
by the redaction tests of Chapter~\ref{cap:validation}. The CV and
job-description parsers additionally write per-call telemetry rows, giving a
queryable view of provider reliability, latency and token cost.

%----------------------------------------------------------------------------------------

\section{Continuous Integration and Deployment}
\label{sec:impl-cicd}

Every change runs through GitHub Actions, which type-checks the codebase with
\texttt{tsc -{}-noEmit} and executes the full Vitest suite with coverage,
uploading the coverage report as an artefact. The canonical local validation
mirrors the first step, since a full production build depends on
deployment-only environment variables. Deployment is handled by Vercel, which
builds with Turbopack and provisions a preview deployment per change, so the
live deployment is always the demonstrable product increment referred to in
Chapter~\ref{cap:analysis}. A repository policy (\texttt{legacy-peer-deps})
keeps installs reproducible by pinning peers explicitly where they are only
transitively required.

%----------------------------------------------------------------------------------------

\section{Selected Code Excerpts}
\label{sec:impl-code}

Three short excerpts illustrate the patterns described above. The first is
the core of the matching engine: a pure function that scores only the
applicable, weighted criteria and computes eligibility as their conjunction.

\begin{lstlisting}[language=TypeScript,caption={The matching engine core (\texttt{computeJobFit}).},captionpos=b]
export function computeJobFit(job, candidate, config = DEFAULT_FIT_CONFIG) {
  const breakdown = [
    matchField(job, candidate),    matchExperience(job, candidate),
    matchSeniority(job, candidate), matchRequiredSkills(job, candidate, config),
    matchPreferredSkills(job, candidate), matchLanguages(job, candidate),
    matchEducation(job, candidate),
  ];
  const weights = config.criterionWeights ?? DEFAULT_CRITERION_WEIGHTS;
  // Only applicable criteria with a positive weight contribute.
  const contributing = breakdown.filter(
    (c) => c.applicable && (weights[c.key] ?? 0) > 0);
  const totalWeight = contributing.reduce((s, c) => s + (weights[c.key] ?? 0), 0);
  const overallScore = totalWeight === 0 ? 0 : Math.round(
    contributing.reduce((s, c) => s + c.score * (weights[c.key] ?? 0), 0)
      / totalWeight);
  const isEligible = contributing.every((c) => c.met);
  return { overallScore, isEligible, breakdown };
}
\end{lstlisting}

The second is the heart of the authorisation middleware: a single guard that
gates every API route, returning \texttt{401} for unauthenticated callers and
\texttt{403} for non-HR callers on the HR-only prefixes.

\begin{lstlisting}[language=TypeScript,caption={Centralised API route gating in the middleware.},captionpos=b]
const role = user?.app_metadata?.role as "hr" | "candidate" | undefined;
if (pathname.startsWith("/api/")) {
  const isPublic = PUBLIC_API_PREFIXES.some((p) => pathname.startsWith(p));
  if (!isPublic) {
    if (!user)
      return NextResponse.json({ error: "Unauthorized" }, { status: 401 });
    const hrOnly = HR_ONLY_API_PREFIXES.some((p) => pathname.startsWith(p));
    if (hrOnly && role !== "hr")
      return NextResponse.json({ error: "Forbidden" }, { status: 403 });
  }
}
\end{lstlisting}

The third illustrates the data-integrity boundary: untrusted search input is
escaped before it is interpolated into a PostgREST \texttt{.or()} filter, so
special characters cannot alter the query.

\begin{lstlisting}[language=TypeScript,caption={Escaping user input before a PostgREST filter.},captionpos=b]
const term = escapeOrTerm(rawSearch);
query = query.or(`name.ilike.%${term}%,email.ilike.%${term}%`);
\end{lstlisting}
